% --- Archon physics formalization source begin ---
% archon:physics
% archon:covers PhyXMiniProblems/problem_phyx_mini_0059.lean
% archon:source-report reports/phyx_mini/problem_phyx_mini_0059.source.json

\chapter{Physics problem phyx\_mini\_0059}
\label{ch:PhyXMiniProblems_problem_phyx_mini_0059}

\paragraph{Problem source.}
\#\# Physical scenario

A 3.0-cm-high object is located 20 cm from a concave mirror. The mirror’s radius of curvature is 80 cm.

\#\# Question

Determine the height of the image.

\#\# Answer choices

- A: A: 1.4cm

- B: B: 5.6cm

- C: C: 6.0cm

- D: D: 2.4cm

\#\# Figure caption (auxiliary; use the image as primary evidence)

The image is a diagram illustrating the path of light rays in a concave mirror setup. 

Key elements include:

1. **Mirror**: A concave mirror is depicted with its principal axis running horizontally through the center. The mirror has a reflective surface on the left side.

2. **Object**: A green arrow represents the object situated 20 cm from the mirror on the principal axis.

3. **Focal Point**: Marked as \textbackslash{}( f = 40 \textbackslash{}, \textbackslash{}text\{cm\} \textbackslash{}), indicating the focal length of the mirror.

4. **Image**: A pink arrow on the right side of the mirror, labeled "Virtual image," shows the position and orientation of the virtual image formed. It appears to be located behind the mirror.

5. **Light Rays**: Multiple purple lines demonstrate the path of light rays emanating from the object:
   - Solid lines show rays reflecting off the mirror.
   - Dashed lines trace the extensions of these rays meeting at the point of the virtual image.

6. **Distances**: 
   - The object distance from the mirror is marked as \textbackslash{}( s = 20 \textbackslash{}, \textbackslash{}text\{cm\} \textbackslash{}).
   - The distance to the virtual image is indicated by \textbackslash{}( s' \textbackslash{}).

7. **Mirror Plane**: A dotted line perpendicular to the principal axis marks the position of the mirror plane.

This setup is a classic example in optics to show how a concave mirror forms a virtual image when the object is placed inside the focal length.

\#\# Dataset metadata

Geometrical Optics; Physical Model Grounding Reasoning, Spatial Relation Reasoning

\paragraph{Recorded answer/context.}
C: C: 6.0cm

\paragraph{Figure/image path.}
/root/proposal\_for\_physic/science-mango/phyx\_mini\_run/phyx\_data/test\_image/59.png

\paragraph{Formalization target.}
create a compiling Lean file with sorry bodies at `PhyXMiniProblems/problem\_phyx\_mini\_0059.lean`. The Lean declarations must preserve the physical quantities, dimensions or dimensional roles, figure labels, governing-law hypotheses, and final relation expressed by this problem.
Use Mathlib/Physlib names found through LeanExplore where available. If a physics API is missing, introduce faithful local abstractions rather than scalar placeholder aliases.

\begin{theorem}[Physics formalization target]
\label{thm:physics:phyx_mini_0059:target}
\lean{PhyXMiniProblems.ProblemPhyXMini0059.problem_phyx_mini_0059}
\uses{def:physics:phyx-mini-0059:phyxminiproblems-problemphyxmini0059-lengthincentimeters, def:physics:phyx-mini-0059:phyxminiproblems-problemphyxmini0059-concavemirrorsetup, def:physics:phyx-mini-0059:phyxminiproblems-problemphyxmini0059-matchesfigurereadouts, def:physics:phyx-mini-0059:phyxminiproblems-problemphyxmini0059-satisfieslabeledaxialgeometry, def:physics:phyx-mini-0059:phyxminiproblems-problemphyxmini0059-usescartesianmirrorsignconvention, def:physics:phyx-mini-0059:phyxminiproblems-problemphyxmini0059-satisfiessphericalmirrorfocallaw, def:physics:phyx-mini-0059:phyxminiproblems-problemphyxmini0059-satisfiesgaussianmirrorequation, def:physics:phyx-mini-0059:phyxminiproblems-problemphyxmini0059-satisfiessignedmagnificationlaw, def:physics:phyx-mini-0059:phyxminiproblems-problemphyxmini0059-satisfiesimageheightlaw, def:physics:phyx-mini-0059:phyxminiproblems-problemphyxmini0059-matchesanswer}
The assigned autoformalize agent should translate this physics problem into Lean declarations in the covered file, with theorem and lemma proof bodies written as `by sorry`.
\end{theorem}
\begin{proof}
This is an autoformalization task, not a proof task. Produce statements that can later be proved by the physics prover without weakening the physical model.
\end{proof}

% --- Archon declaration topology begin ---
\section{Declaration topology}
\begin{definition}[Length Quantity]
  \label{def:physics:phyx-mini-0059:phyxminiproblems-problemphyxmini0059-lengthquantity}
  \lean{PhyXMiniProblems.ProblemPhyXMini0059.LengthQuantity}
  A signed physical length, independent of the unit chosen to read it
\end{definition}
\begin{proof}
  This is the declared model definition of Length Quantity.
\end{proof}
\begin{definition}[centimeter Unit Choices]
  \label{def:physics:phyx-mini-0059:phyxminiproblems-problemphyxmini0059-centimeterunitchoices}
  \lean{PhyXMiniProblems.ProblemPhyXMini0059.centimeterUnitChoices}
  Unit choices with centimeters selected as the unit of length
\end{definition}
\begin{proof}
  This is the declared model definition of centimeter Unit Choices.
\end{proof}
\begin{definition}[length In Centimeters]
  \label{def:physics:phyx-mini-0059:phyxminiproblems-problemphyxmini0059-lengthincentimeters}
  \lean{PhyXMiniProblems.ProblemPhyXMini0059.lengthInCentimeters}
  \uses{def:physics:phyx-mini-0059:phyxminiproblems-problemphyxmini0059-lengthquantity, def:physics:phyx-mini-0059:phyxminiproblems-problemphyxmini0059-centimeterunitchoices}
  The signed scalar readout of a physical length in centimeters
\end{definition}
\begin{proof}
  This is the declared model definition of length In Centimeters.
\end{proof}
\begin{definition}[Axial Figure Point]
  \label{def:physics:phyx-mini-0059:phyxminiproblems-problemphyxmini0059-axialfigurepoint}
  \lean{PhyXMiniProblems.ProblemPhyXMini0059.AxialFigurePoint}
  Labeled points on the horizontal principal axis in the primary figure
\end{definition}
\begin{proof}
  This is the declared model definition of Axial Figure Point.
\end{proof}
\begin{definition}[Spherical Mirror Kind]
  \label{def:physics:phyx-mini-0059:phyxminiproblems-problemphyxmini0059-sphericalmirrorkind}
  \lean{PhyXMiniProblems.ProblemPhyXMini0059.SphericalMirrorKind}
  The spherical-mirror type as viewed from the incident-light side
\end{definition}
\begin{proof}
  This is the declared model definition of Spherical Mirror Kind.
\end{proof}
\begin{definition}[Image Nature]
  \label{def:physics:phyx-mini-0059:phyxminiproblems-problemphyxmini0059-imagenature}
  \lean{PhyXMiniProblems.ProblemPhyXMini0059.ImageNature}
  Whether reflected rays themselves, or only their backward extensions, meet
\end{definition}
\begin{proof}
  This is the declared model definition of Image Nature.
\end{proof}
\begin{definition}[Image Orientation]
  \label{def:physics:phyx-mini-0059:phyxminiproblems-problemphyxmini0059-imageorientation}
  \lean{PhyXMiniProblems.ProblemPhyXMini0059.ImageOrientation}
  The transverse orientation of the image relative to the object
\end{definition}
\begin{proof}
  This is the declared model definition of Image Orientation.
\end{proof}
\begin{definition}[Ray Construction]
  \label{def:physics:phyx-mini-0059:phyxminiproblems-problemphyxmini0059-rayconstruction}
  \lean{PhyXMiniProblems.ProblemPhyXMini0059.RayConstruction}
  How the reflected rays in the primary figure locate the image
\end{definition}
\begin{proof}
  This is the declared model definition of Ray Construction.
\end{proof}
\begin{definition}[Mirror Plane Relation]
  \label{def:physics:phyx-mini-0059:phyxminiproblems-problemphyxmini0059-mirrorplanerelation}
  \lean{PhyXMiniProblems.ProblemPhyXMini0059.MirrorPlaneRelation}
  The relation drawn between the mirror plane and principal axis
\end{definition}
\begin{proof}
  This is the declared model definition of Mirror Plane Relation.
\end{proof}
\begin{definition}[Concave Mirror Setup]
  \label{def:physics:phyx-mini-0059:phyxminiproblems-problemphyxmini0059-concavemirrorsetup}
  \lean{PhyXMiniProblems.ProblemPhyXMini0059.ConcaveMirrorSetup}
  \uses{def:physics:phyx-mini-0059:phyxminiproblems-problemphyxmini0059-lengthquantity, def:physics:phyx-mini-0059:phyxminiproblems-problemphyxmini0059-axialfigurepoint, def:physics:phyx-mini-0059:phyxminiproblems-problemphyxmini0059-sphericalmirrorkind, def:physics:phyx-mini-0059:phyxminiproblems-problemphyxmini0059-imagenature, def:physics:phyx-mini-0059:phyxminiproblems-problemphyxmini0059-imageorientation, def:physics:phyx-mini-0059:phyxminiproblems-problemphyxmini0059-rayconstruction, def:physics:phyx-mini-0059:phyxminiproblems-problemphyxmini0059-mirrorplanerelation}
  The physical quantities and labeled axial positions in the mirror diagram. objectDistance is the positive figure label s. The field signedImageDistance is the signed Gaussian distance s': it is negative for the pictured virtual image. The radius field records the positive magnitude given in the problem
\end{definition}
\begin{proof}
  This is the declared model definition of Concave Mirror Setup.
\end{proof}
\begin{definition}[axial Separation Centimeters]
  \label{def:physics:phyx-mini-0059:phyxminiproblems-problemphyxmini0059-axialseparationcentimeters}
  \lean{PhyXMiniProblems.ProblemPhyXMini0059.axialSeparationCentimeters}
  \uses{def:physics:phyx-mini-0059:phyxminiproblems-problemphyxmini0059-lengthincentimeters, def:physics:phyx-mini-0059:phyxminiproblems-problemphyxmini0059-axialfigurepoint, def:physics:phyx-mini-0059:phyxminiproblems-problemphyxmini0059-concavemirrorsetup}
  The directed centimeter separation of two labels on the principal axis
\end{definition}
\begin{proof}
  This is the declared model definition of axial Separation Centimeters.
\end{proof}
\begin{definition}[Matches Figure Readouts]
  \label{def:physics:phyx-mini-0059:phyxminiproblems-problemphyxmini0059-matchesfigurereadouts}
  \lean{PhyXMiniProblems.ProblemPhyXMini0059.MatchesFigureReadouts}
  \uses{def:physics:phyx-mini-0059:phyxminiproblems-problemphyxmini0059-lengthincentimeters, def:physics:phyx-mini-0059:phyxminiproblems-problemphyxmini0059-concavemirrorsetup}
  Problem-statement and primary-figure readouts. These include only supplied data: the object dimensions and placement, the radius and displayed focal length, and the qualitative nature and orientation of the pictured image. The requested image height is deliberately absent
\end{definition}
\begin{proof}
  This is the declared model definition of Matches Figure Readouts.
\end{proof}
\begin{definition}[Satisfies Labeled Axial Geometry]
  \label{def:physics:phyx-mini-0059:phyxminiproblems-problemphyxmini0059-satisfieslabeledaxialgeometry}
  \lean{PhyXMiniProblems.ProblemPhyXMini0059.SatisfiesLabeledAxialGeometry}
  \uses{def:physics:phyx-mini-0059:phyxminiproblems-problemphyxmini0059-lengthincentimeters, def:physics:phyx-mini-0059:phyxminiproblems-problemphyxmini0059-concavemirrorsetup, def:physics:phyx-mini-0059:phyxminiproblems-problemphyxmini0059-axialseparationcentimeters}
  The labeled separations agree with the physical distance fields. Axial coordinates increase to the right, so a virtual image behind the mirror has signed Gaussian distance equal to the negative mirror-to-image separation
\end{definition}
\begin{proof}
  This is the declared model definition of Satisfies Labeled Axial Geometry.
\end{proof}
\begin{definition}[Uses Cartesian Mirror Sign Convention]
  \label{def:physics:phyx-mini-0059:phyxminiproblems-problemphyxmini0059-usescartesianmirrorsignconvention}
  \lean{PhyXMiniProblems.ProblemPhyXMini0059.UsesCartesianMirrorSignConvention}
  \uses{def:physics:phyx-mini-0059:phyxminiproblems-problemphyxmini0059-lengthincentimeters, def:physics:phyx-mini-0059:phyxminiproblems-problemphyxmini0059-concavemirrorsetup}
  Positivity and signed-distance conventions for the depicted optical branch
\end{definition}
\begin{proof}
  This is the declared model definition of Uses Cartesian Mirror Sign Convention.
\end{proof}
\begin{definition}[Satisfies Spherical Mirror Focal Law]
  \label{def:physics:phyx-mini-0059:phyxminiproblems-problemphyxmini0059-satisfiessphericalmirrorfocallaw}
  \lean{PhyXMiniProblems.ProblemPhyXMini0059.SatisfiesSphericalMirrorFocalLaw}
  \uses{def:physics:phyx-mini-0059:phyxminiproblems-problemphyxmini0059-concavemirrorsetup}
  The focal length of a spherical mirror is half the magnitude of its radius of curvature. Requiring the relation in every unit choice makes the governing law independent of the centimeter readout used by the figure
\end{definition}
\begin{proof}
  This is the declared model definition of Satisfies Spherical Mirror Focal Law.
\end{proof}
\begin{definition}[Satisfies Gaussian Mirror Equation]
  \label{def:physics:phyx-mini-0059:phyxminiproblems-problemphyxmini0059-satisfiesgaussianmirrorequation}
  \lean{PhyXMiniProblems.ProblemPhyXMini0059.SatisfiesGaussianMirrorEquation}
  \uses{def:physics:phyx-mini-0059:phyxminiproblems-problemphyxmini0059-concavemirrorsetup}
  The signed Gaussian mirror equation 1/f = 1/p + 1/q, in the division-free dimensionally homogeneous form f * (p + q) = p * q and in every unit system
\end{definition}
\begin{proof}
  This is the declared model definition of Satisfies Gaussian Mirror Equation.
\end{proof}
\begin{definition}[Satisfies Signed Magnification Law]
  \label{def:physics:phyx-mini-0059:phyxminiproblems-problemphyxmini0059-satisfiessignedmagnificationlaw}
  \lean{PhyXMiniProblems.ProblemPhyXMini0059.SatisfiesSignedMagnificationLaw}
  \uses{def:physics:phyx-mini-0059:phyxminiproblems-problemphyxmini0059-concavemirrorsetup}
  The signed transverse-magnification law m = -q/p, written without division. The magnification is dimensionless, while p and q are signed lengths
\end{definition}
\begin{proof}
  This is the declared model definition of Satisfies Signed Magnification Law.
\end{proof}
\begin{definition}[Satisfies Image Height Law]
  \label{def:physics:phyx-mini-0059:phyxminiproblems-problemphyxmini0059-satisfiesimageheightlaw}
  \lean{PhyXMiniProblems.ProblemPhyXMini0059.SatisfiesImageHeightLaw}
  \uses{def:physics:phyx-mini-0059:phyxminiproblems-problemphyxmini0059-concavemirrorsetup}
  The image-height law hᵢ = m hₒ, required in every unit system
\end{definition}
\begin{proof}
  This is the declared model definition of Satisfies Image Height Law.
\end{proof}
\begin{definition}[Answer Choice]
  \label{def:physics:phyx-mini-0059:phyxminiproblems-problemphyxmini0059-answerchoice}
  \lean{PhyXMiniProblems.ProblemPhyXMini0059.AnswerChoice}
  Labels of the four multiple-choice answers
\end{definition}
\begin{proof}
  This is the declared model definition of Answer Choice.
\end{proof}
\begin{definition}[answer Height In Centimeters]
  \label{def:physics:phyx-mini-0059:phyxminiproblems-problemphyxmini0059-answerheightincentimeters}
  \lean{PhyXMiniProblems.ProblemPhyXMini0059.answerHeightInCentimeters}
  \uses{def:physics:phyx-mini-0059:phyxminiproblems-problemphyxmini0059-answerchoice}
  The image height in centimeters printed beside each answer choice
\end{definition}
\begin{proof}
  This is the declared model definition of answer Height In Centimeters.
\end{proof}
\begin{definition}[Matches Answer]
  \label{def:physics:phyx-mini-0059:phyxminiproblems-problemphyxmini0059-matchesanswer}
  \lean{PhyXMiniProblems.ProblemPhyXMini0059.MatchesAnswer}
  \uses{def:physics:phyx-mini-0059:phyxminiproblems-problemphyxmini0059-lengthquantity, def:physics:phyx-mini-0059:phyxminiproblems-problemphyxmini0059-lengthincentimeters, def:physics:phyx-mini-0059:phyxminiproblems-problemphyxmini0059-answerchoice, def:physics:phyx-mini-0059:phyxminiproblems-problemphyxmini0059-answerheightincentimeters}
  Exact agreement between a physical image height and a displayed choice
\end{definition}
\begin{proof}
  This is the declared model definition of Matches Answer.
\end{proof}
\begin{lemma}[signed image distance and magnification]
  \label{lem:physics:phyx-mini-0059:phyxminiproblems-problemphyxmini0059-signed-image-distance-and-magnification}
  \lean{PhyXMiniProblems.ProblemPhyXMini0059.signed_image_distance_and_magnification}
  \uses{def:physics:phyx-mini-0059:phyxminiproblems-problemphyxmini0059-lengthincentimeters, def:physics:phyx-mini-0059:phyxminiproblems-problemphyxmini0059-concavemirrorsetup, def:physics:phyx-mini-0059:phyxminiproblems-problemphyxmini0059-matchesfigurereadouts, def:physics:phyx-mini-0059:phyxminiproblems-problemphyxmini0059-satisfieslabeledaxialgeometry, def:physics:phyx-mini-0059:phyxminiproblems-problemphyxmini0059-usescartesianmirrorsignconvention, def:physics:phyx-mini-0059:phyxminiproblems-problemphyxmini0059-satisfiessphericalmirrorfocallaw, def:physics:phyx-mini-0059:phyxminiproblems-problemphyxmini0059-satisfiesgaussianmirrorequation, def:physics:phyx-mini-0059:phyxminiproblems-problemphyxmini0059-satisfiessignedmagnificationlaw}
  The governing mirror and magnification equations determine the unprinted intermediate quantities: the virtual image distance is -40 cm and the upright transverse magnification is 2
\end{lemma}
\begin{proof}
  The conclusion follows from the governing relations and figure readouts named in the statement of signed image distance and magnification.
\end{proof}
% --- Archon declaration topology end ---
% --- Archon physics formalization source end ---
